\documentclass[12pt]{article}
\usepackage[utf8]{inputenc}
\usepackage[T1]{fontenc}
\usepackage[margin=1in]{geometry}
\usepackage{setspace}
\usepackage{xcolor}
\usepackage[colorlinks=true, urlcolor=blue!60!black, linkcolor=blue!60!black]{hyperref}
\setlength{\parindent}{0em}
\setlength{\parskip}{0.6em}
\renewcommand{\baselinestretch}{1.0}
\begin{document}
\begin{center}
{\LARGE \textit{Nixon v. Administrator of General Services}, 433 U.S. 425 (1977)\par}
\vspace{0.5em}
{\large Case Brief --- Two-Page Edition\par}
\end{center}

\section{Memory Jogger}
Congress could seize and screen Nixon's presidential papers and tapes: no bill of attainder (a ``legitimate class of one,'' not punished) and no separation-of-powers violation (\href{https://supreme.justia.com/cases/federal/us/433/425/}{Justia}).

\section{Detailed Case Facts}
After Nixon resigned, a GSA agreement let him control his 42 million pages and 880 tapes and destroy the tapes. Congress abrogated it in the Presidential Recordings and Materials Preservation Act, ordering GSA custody, archival screening, return of personal materials, and compensation for any taking.

\section{Procedural History}
A three-judge district court rejected Nixon's facial challenge; the Court noted probable jurisdiction.

\section{Judicial Votes}
Affirmed 7--2. Brennan wrote for the Court; Stevens, White, Blackmun, and Powell concurred in whole or part; Burger and Rehnquist dissented.

\section{Holding}
Nixon was a legitimate class of one, and the Act inflicted no punishment under historical, functional, or motivational tests. It did not unduly prevent the Executive from performing its functions, and his privilege, privacy, and associational claims yielded to minimal archival screening.

\section{Analysis of Opinions}
\textbf{Majority:} only Nixon's papers faced destruction; compensation negated confiscation; the Act served nonpunitive purposes. Separation of powers does not require ``three air-tight departments''; custody stayed with executive archivists, and two Presidents supported the law. \textbf{Stevens:} Nixon's resignation and pardon made him unique. \textbf{Burger:} the Act was an attainder. \textbf{Rehnquist:} it is a ``veritable sword of Damocles'' over future Presidents.

\section{Examples: Future Applications}
\textbf{Same-Side:} a statute securing one former official's email server, with screening and compensation. \textbf{Opposite-Side:} uncompensated forfeiture of a named company's patents amid condemnation (\textit{Foretich}). \textbf{Fence-Sitter:} a hostile successor's waiver of privilege for publication.

\section{Critique}
The majority turned the Bill of Attainder Clause into a deferential inquiry, risked normalizing pretextual ``classes of one,'' and leaned on incumbent consent that may fail in polarized transitions.

\section{Key Quotations}
\begin{quote}
``The Court, therefore, often has looked beyond mere historical experience and has applied a functional test of the existence of punishment, analyzing whether the law under challenge, viewed in terms of the type and severity of burdens imposed, reasonably can be said to further nonpunitive legislative purposes.''
\end{quote}
\noindent --- Justice Brennan, for the Court, 433 U.S. at 475--476.

\end{document}
